\section{Введение}

AXI4-Lite Cache --- параметризуемый наборно-ассоциативный кэш, снижающий среднюю задержку доступа процессора к внешней памяти.

\subsection{Характеристики}

\begin{table}[htbp]
	\centering
	\label{tab}
	\small
	\begin{tabularx}{\textwidth}{l X}
		\toprule
		\textbf{Характеристика} & \textbf{Описание} \\
		\midrule
		
		\textit{Степень ассоциативности} &
		Настраиваемое число секций кэша. \\
		
		\textit{Разрядность адреса и данных} &
		Настраиваемая разрядность адреса и данных. \\
		
		\textit{Число наборов} &
		Настраиваемое число наборов кэша. \\
		
		\textit{Write through} &
		Запись немедленно передаётся во внешнюю память. \\
		
		\textit{Стратегия замещения NMRU} &
		Индекс секции-жерты выбирается случайно, наиболее недавно использованная секция исключается из замещения. \\
		
		\textit{Размер строки кэша} &
		Размер строки кэша равен разрядности данных. \\
		
		\textit{Интерфейс AXI4-Lite} &
		Интерфейсы AXI4-Lite со стороны процессора и внешней памяти. \\
		
		\textit{Нет outstanding транзакций} &
		Кэш обрабатывает одну транзакцию за раз. \\
		
		\textit{Статистика попаданий и промахов} &
		Подсчёт попаданий и промахов только в симуляции. \\

		\textit{Латентность} &
		1 такт при попадании в кэш. \\
		
		\bottomrule
	\end{tabularx}
	\caption{Характеристики кэша}
\end{table}

% --- РАЗДЕЛ 2 ---
